% ============================================================
%  BEISPIEL-KAPITEL — zeigt alle Makros in Verwendung
%  Als parts/01_kapitel.tex speichern und in TEMPLATE.tex einbinden
% ============================================================

% ---- \bereich[Farbe]{Titel} --------------------------------
% Farbe muss in TEMPLATE.tex mit \colorlet definiert sein
\bereich[cKap1]{Abschnittstitel}

% ---- formelreihe -------------------------------------------
% 3 Spalten, getrennt mit &, Zeile endet mit \\
% Empfehlung: Spalte 1+2 = Formeln, Spalte 3 = Beispiele
\begin{formelreihe}
  % --- Zeile 1: normaler Fall ---
  \formelblock{Formelname}{f(x) = \int_a^b g(x)\,\mathrm{d}x} &
  \formelblock{Weitere Formel}{P(A) = \dfrac{|A|}{|\Omega|}} &
  \beispielblock{Kontext}{P(\text{Kopf}) = \dfrac{1}{2}} \\[0.3em]

  % --- Zeile 2: mehrzeilige Formel mit aligned ---
  \formelblock{Mehrzeilig}{\begin{aligned}
    \mu &= E\{X\} \\
    \sigma^2 &= E\{X^2\} - \mu^2
  \end{aligned}} &
  % --- leere Zelle ---
  &
  % --- Beispiel mit Zeilenumbruch (aligned in #2) ---
  \beispielblock{Würfel}{%
    E\{X\}=\sum_{k=1}^{6}k\cdot\tfrac{1}{6}=3{.}5 \\[0.15em]
    \text{(aligned nicht nötig für 1 Zeile)}} \\[0.3em]

  % --- Zeile 3: Infoblock in Spalte 1 ---
  \infoblock{Erklärungstext erscheint\\kursiv und klein.\\
             Gut für Definitionen.} &
  \formelblock{Formel in Sp. 2}{\sigma_X = \sqrt{\sigma_X^2}} &
  \\
\end{formelreihe}
\bereichende

% ---- Zweiter Abschnitt -------------------------------------
\bereich[cKap1]{Zweiter Abschnitt}
\begin{formelreihe}
  \formelblock{Beispiel}{y = ax + b} &
  \formelblock{Folge}{\mu_y = a\mu_x + b,\quad \sigma_y^2 = a^2\sigma_x^2} &
  \\
\end{formelreihe}
\bereichende
